\documentclass[10pt,a4paper]{amsart}
\usepackage[margin=19mm]{geometry}
\usepackage{amsmath,amssymb,mathtools}
\usepackage[scr=rsfs,cal=euler]{mathalfa}
\usepackage{xfrac}
\usepackage[unicode,hidelinks,bookmarks=false]{hyperref}
\newcommand{\ie}{i.e.\ }
\newcommand{\cf}{cf.\ }
\newcommand{\resp}{respectively\ }
\newcommand{\Subsection}{\subsection}
\providecommand{\nobreakdash}{\nobreak\hbox{-}}
\setlength{\emergencystretch}{2em}
\multlinegap0pt
\usepackage{mathtools}
\usepackage{amssymb}
\usepackage[shortlabels]{enumitem}
\usepackage{csquotes}
\newtheorem{theorem}{Theorem}
\newcommand{\Aut}{\operatorname{Aut}}
\newcommand{\Out}{\operatorname{Out}}
\newcommand{\Stab}{\operatorname{Stab}}
\newcommand{\ad}{\operatorname{ad}}
\newcommand{\cE}{\mathcal{E}}
\newcommand{\cP}{\mathcal{P}}
\newcommand{\cusp}{\operatorname{cusp}}
\title{Pinned Jordan Decomposition of Characters and Depth-Zero Hecke Algebras}
\author{Prashant Arote, Manish Mishra}
\date{Source: arXiv:2605.03036v2 (CC BY 4.0)}
\begin{document}
\maketitle

\noindent\emph{Source and licence.} Pinned Jordan Decomposition of Characters and Depth-Zero Hecke Algebras. Version arXiv:2605.03036v2 (CC BY 4.0), distributed by arXiv under the Creative Commons Attribution 4.0 International licence, http://creativecommons.org/licenses/by/4.0/, as recorded in the arXiv licence metadata for this exact version and shown on the abstract page https://arxiv.org/abs/2605.03036v2. Full licence notice: \texttt{licenses/arxiv-2605.03036-CC-BY-4.0.txt}.\par\smallskip

\noindent\textit{Editorial scope.} Counted unit: the article's theorem thm:pinned-Malle-matching-simple (source0.tex lines 3431--3482). The article's complete original proof of it is delivered with it (source0.tex lines 3483--3635, one \texttt{proof} environment of 153 lines). No mathematics is written, repaired or re-derived: the statement and the proof are the article's own lines, in the article's own notation, and no retained line is substituted or re-flowed.  No further source-local context is needed: every cross-reference of every retained line resolves inside the two delivered spans.  Nothing was omitted from any retained span, so the package carries no omission record.  No retained line contains a citation, so the wrapper carries no bibliography and no bibliography entry.  No retained line contains a figure environment, an image-inclusion command or drawing code, so no image is delivered and the package declares no asset dependency.  Editorial formatting only, in addition: an AMS \texttt{amsart} preamble that restates the article's own declarations and macros used by the retained lines, namely the environments \texttt{theorem} on separate counters, together with the standard packages those lines need.  The retained environments print in this document as Theorem 1 in order of appearance, whereas the article prints the same items with the numbering of its own full document.  The complete native source of the article as distributed in the arXiv e-print for this exact version is bundled unchanged as \texttt{sources/199763/source0.tex}.
\begin{theorem}[Pinned Malle matching for simple groups]
\label{thm:pinned-Malle-matching-simple}
Let \(\mathbf G\) be simple and simply connected, equipped with a Frobenius
endomorphism \(F\), and put \(G=\mathbf G^F\).  Fix an \(F\)-stable pinning
\(\cP\) of \(\mathbf G\), and let \(\psi_{\cP}\) be the corresponding
non-degenerate character of the unipotent radical of the fixed Borel subgroup.
Let
\[
        A:=G_{\ad}/G
\]
denote the group of diagonal automorphisms of \(G\).  Let
\(s\in G^{*F}\), and let
\[
        \mathcal S_s\subset \cE(G,s)
\]
be the \(A\)-orbit of semisimple characters in \(\cE(G,s)\).  Denote by
\[
        \rho_{s,\psi_{\cP}}\in \mathcal S_s
\]
the Whittaker-normalized semisimple character determined by \(\psi_{\cP}\).

Let
\[
        F_O\subset \cE(G,s)_{\cusp}
\]
be an \(A\)-orbit of cuspidal characters.  Then the pinning \(\cP\) determines
a distinguished element
\[
        \rho_{O,\cP}\in F_O
\]
such that
\[
        \Aut(G)_{\rho_{O,\cP}}
        =
        \Aut(G)_{\rho_{s,\psi_{\cP}}}.
\]
Consequently there is a unique \(A\)-equivariant bijection
\[
        m_{\cP,O}:F_O\xrightarrow{\sim}\mathcal S_s
\]
satisfying
\[
        m_{\cP,O}(\rho_{O,\cP})=\rho_{s,\psi_{\cP}}.
\]
Moreover, for every \(\rho\in F_O\), one has
\[
        \Aut(G)_\rho
        =
        \Aut(G)_{m_{\cP,O}(\rho)}.
\]
We call \(m_{\cP,O}\) the \(\cP\)-preferred Malle matching on the orbit \(F_O\).
\end{theorem}

\begin{proof}
Write the action of \(a\in A\) on characters as
\[
        \eta^a:=\eta\circ \ad(a).
\]
Malle's theorem gives, for each cuspidal character
\(\rho\in \cE(G,s)\), a semisimple character
\(\chi\in \cE(G,s)\) satisfying
\[
        \Aut(G)_\chi=\Aut(G)_\rho.
\]
In particular, the \(A\)-orbit \(F_O\) is isomorphic, as an \(A\)-set, to the
\(A\)-orbit \(\mathcal S_s\) of semisimple characters.  The point is to choose
the isomorphism in a way normalized by the pinning.

We first treat type \(A\).  In this case the diagonal automorphism group may
have order larger than two, so the stabilizer condition alone does not choose
a distinguished point.  Instead one uses the Whittaker normalization.  By the
type \(A\) equivariant character correspondence of Cabanes--Späth, as used in
Malle's proof, each diagonal orbit of cuspidal characters in \(\cE(G,s)\)
contains a unique \(\psi_{\cP}\)-generic constituent.  We define
\[
        \rho_{O,\cP}
\]
to be this constituent.  The semisimple orbit \(\mathcal S_s\) similarly
contains the unique \(\psi_{\cP}\)-generic semisimple character
\(\rho_{s,\psi_{\cP}}\).  The type \(A\) equivariance theorem gives
\[
        \Aut(G)_{\rho_{O,\cP}}
        =
        \Aut(G)_{\rho_{s,\psi_{\cP}}}.
\]

We now assume that \(\mathbf G\) is not of type \(A\).  Let
\[
        S_\rho:=\Aut(G)_\rho,\qquad A_\rho:=S_\rho\cap A,
\]
and let \(B_\rho\) be the image of \(S_\rho\) in the graph-field part of
\(\Out(G)\).  If \(\chi\in\mathcal S_s\) is one semisimple character satisfying
\[
        \Aut(G)_\chi=S_\rho,
\]
then the other diagonal translates of \(\chi\) with the same stabilizer
\(S_\rho\) are parametrized by
\[
        (A/A_\rho)^{B_\rho}.
\]
For the non-type \(A\) groups occurring in Malle's reduction, this group is
either trivial or of order two.

If
\[
        (A/A_\rho)^{B_\rho}=1,
\]
then there is no residual choice: the semisimple character in
\(\mathcal S_s\) whose full automorphism stabilizer equals \(S_\rho\) is unique.
Equivalently, there is a unique element of \(F_O\) whose stabilizer equals
\(\Aut(G)_{\rho_{s,\psi_{\cP}}}\).  We define this element to be
\(\rho_{O,\cP}\).

It remains to discuss the residual twofold cases.  These are precisely the
cases in Malle's proof where the stabilizer condition leaves two diagonal
translates undistinguished.  Malle's proof treats them by explicit Brauer-tree
chains, by Deligne--Lusztig induction chains with multiplicity-one
constituents, and by the dihedral actions occurring in type \(D\).  The pinning
removes the remaining sign ambiguity as follows.

In the Brauer-tree cases, Malle connects the cuspidal character to the
semisimple character by a labelled chain of non-exceptional vertices on
Brauer trees.  We root this chain at the pinned semisimple endpoint
\(\rho_{s,\psi_{\cP}}\).  Since the labels of hooks, cohooks, and the relevant
simple factors are fixed by the pinned root datum, this rooted chain has a
unique cuspidal endpoint in \(F_O\).  We define this endpoint to be
\(\rho_{O,\cP}\).  Malle's Brauer-tree argument then gives equality of the
full automorphism stabilizers.

In the Deligne--Lusztig induction cases, choose the Levi subgroup \(L\) to be
the standard Levi determined by the pinned root datum.  The character on \(L\)
appearing in Malle's induction chain is chosen with the corresponding pinned
Whittaker normalization.  The relevant constituents of
\[
        R_L^G(\psi)
\]
occur with multiplicity one.  Hence, among the two possible diagonal translates
in the residual orbit, the pinned induction datum selects a unique constituent.
This constituent is defined to be \(\rho_{O,\cP}\), and Malle's
Deligne--Lusztig argument gives
\[
        \Aut(G)_{\rho_{O,\cP}}
        =
        \Aut(G)_{\rho_{s,\psi_{\cP}}}.
\]

Finally, in the type \(D\) cases where the residual ambiguity is spinorial, the
non-trivial ambiguity is the interchange of the two half-spin labels, or its
graph-field analogue.  The pinning labels the two terminal vertices of the
\(D_n\)-diagram.  We choose the constituent corresponding to the preferred
terminal vertex determined by \(\cP\).  Malle's case analysis shows that the
cuspidal and semisimple diagonal orbits carry the same dihedral action of the
relevant graph-field automorphisms.  Therefore this pinned choice again has
the same full automorphism stabilizer as the pinned semisimple character.

Thus in every case we have constructed a distinguished element
\[
        \rho_{O,\cP}\in F_O
\]
with
\[
        \Aut(G)_{\rho_{O,\cP}}
        =
        \Aut(G)_{\rho_{s,\psi_{\cP}}}.
\]
In particular,
\[
        \Stab_A(\rho_{O,\cP})
        =
        \Stab_A(\rho_{s,\psi_{\cP}}).
\]
There is therefore a unique \(A\)-equivariant bijection
\[
        m_{\cP,O}:F_O\longrightarrow \mathcal S_s
\]
sending \(\rho_{O,\cP}\) to \(\rho_{s,\psi_{\cP}}\), namely
\[
        m_{\cP,O}\bigl(\rho_{O,\cP}^a\bigr)
        :=
        \rho_{s,\psi_{\cP}}^a
        \qquad(a\in A).
\]
This formula is well-defined because the two basepoints have the same
\(A\)-stabilizer.  It is visibly \(A\)-equivariant and is the only
\(A\)-equivariant map with the prescribed value at \(\rho_{O,\cP}\).

For any \(a\in A\), conjugating the equality
\[
        \Aut(G)_{\rho_{O,\cP}}
        =
        \Aut(G)_{\rho_{s,\psi_{\cP}}}
\]
by \(a\) gives
\[
        \Aut(G)_{\rho_{O,\cP}^a}
        =
        \Aut(G)_{\rho_{s,\psi_{\cP}}^a}.
\]
Hence
\[
        \Aut(G)_\rho
        =
        \Aut(G)_{m_{\cP,O}(\rho)}
\]
for all \(\rho\in F_O\), as claimed.
\end{proof}

\end{document}
